% ===================================================================
%  《战争的雷霆——记一种可能的中国历史》
%  Thunder of War: A Chronological Reconstruction of a Possible
%  Chinese History
%  Moscow University Press / St Petersburg Academic Press, 3rd ed. 2024
%  排版：仿苏式/中译社科类古籍体例
% ===================================================================

\usepackage{geometry}
\geometry{
  paperwidth=185mm, paperheight=260mm,
  inner=21mm, outer=17mm, top=22mm, bottom=24mm,
  headheight=24pt, headsep=6.5mm, footskip=11mm
}

\usepackage{fontspec}
\usepackage{xeCJK}
\usepackage{xcolor}
\usepackage{graphicx}
\usepackage{array,tabularx,booktabs,multirow,makecell}
\usepackage{longtable}
\usepackage{multicol}
\usepackage{calc}
\usepackage{amsmath,amssymb}
\usepackage{enumitem}
\usepackage{fancyhdr}
\usepackage{titlesec}
\usepackage{titletoc}
\usepackage{tcolorbox}
\tcbuselibrary{skins,breakable}
\usepackage[perpage]{footmisc}
\usepackage{etoolbox}
\usepackage{eso-pic}
\usepackage{tikz}
\usetikzlibrary{patterns,arrows.meta,calc,positioning,decorations.pathmorphing,
                decorations.markings,shapes.geometric}
\usepackage{pgfplots}
\pgfplotsset{compat=1.17}
\usepackage{caption}
\usepackage[hidelinks,bookmarksnumbered,bookmarksopen]{hyperref}

% ---------------------------------------------------------------- 颜色
\definecolor{ink}{HTML}{17191C}
\definecolor{graphite}{HTML}{3F454C}
\definecolor{pencil}{HTML}{5B6167}
\definecolor{rulegrey}{HTML}{8B8F93}
\definecolor{faintgrey}{HTML}{C9C9C4}
\definecolor{oxblood}{HTML}{6E1B22}
\definecolor{coverbg}{HTML}{1A1C20}
\definecolor{coverink}{HTML}{E6DFCF}
\definecolor{coverdim}{HTML}{7C7A73}
\definecolor{shade}{HTML}{F2F0EA}
\definecolor{stamp}{HTML}{8C2B2B}

\color{ink}

% ---------------------------------------------------------------- 字体
\setmainfont{NotoSerif}[
  Extension=.ttf,
  UprightFont=*-Regular,
  BoldFont=*-Bold,
  ItalicFont=*-Italic,
  BoldItalicFont=*-BoldItalic,
  Path=/usr/share/fonts/noto/
]
\setsansfont{NotoSans}[
  Extension=.ttf,
  UprightFont=*-Regular,
  BoldFont=*-Bold,
  ItalicFont=*-Italic,
  BoldItalicFont=*-BoldItalic,
  Path=/usr/share/fonts/noto/
]
\setmonofont{JetBrains Maple Mono}[
  UprightFont=* Regular, BoldFont=* Bold, Scale=0.92
]

\setCJKmainfont{Noto Serif CJK SC}[
  BoldFont=Noto Serif CJK SC Bold,
  ItalicFont=Noto Serif CJK SC,
  BoldItalicFont=Noto Serif CJK SC Bold
]
% 让罗马数字、带圈数字等“类汉字”符号使用中文字体
\xeCJKDeclareCharClass{CJK}{"2150 -> "218F, "2460 -> "24FF, "3000 -> "303F}

\setCJKsansfont{Noto Sans CJK SC}[BoldFont=Noto Sans CJK SC Bold]
\setCJKmonofont{Noto Sans Mono CJK SC}

% 伪斜体的字距（中文无真斜体，此处仅作强调标记）
\newcommand{\jia}{\CJKfontspec{Noto Serif CJK SC}}

% ---------------------------------------------------------------- 版心
\linespread{1.42}
\setlength{\parindent}{2em}
\setlength{\parskip}{0pt}
\raggedbottom
\setlength{\emergencystretch}{2em}
\clubpenalty=10000
\widowpenalty=10000

% ---------------------------------------------------------------- 计数
\renewcommand{\thechapter}{\chinese{chapter}}
\renewcommand{\thesection}{\chinese{section}}
\renewcommand{\thesubsection}{\chinese{subsection}}
\renewcommand{\thefigure}{\arabic{chapter}-\arabic{figure}}
\renewcommand{\thetable}{\arabic{chapter}-\arabic{table}}
\renewcommand{\theequation}{\arabic{chapter}-\arabic{equation}}

% ---------------------------------------------------------------- 章标题
\newcommand{\chapterornament}{%
  \vspace{0.8\baselineskip}%
  \begin{tikzpicture}[baseline=0pt]
    \fill[rulegrey] (0,0) rectangle (42mm,0.75pt);
    \fill[rulegrey] (0,2.4pt) rectangle (42mm,2.75pt);
  \end{tikzpicture}}

\titleformat{\chapter}[display]
  {\centering\normalfont}
  {\large\mdseries 第\,\thechapter\,章}
  {0.55\baselineskip}
  {\normalfont\LARGE\bfseries\linespread{1.15}\selectfont}
  [\chapterornament]
\titlespacing*{\chapter}{0pt}{-16pt}{30pt}

\titleformat{\section}
  {\normalfont\sffamily\large\bfseries\linespread{1.2}\selectfont}
  {第\,\thesection\,节}{0.6em}{}
\titlespacing*{\section}{0pt}{18pt}{9pt}

\titleformat{\subsection}
  {\normalfont\normalsize\bfseries}
  {\thesubsection、}{0.2em}{}
\titlespacing*{\subsection}{0pt}{13pt}{6pt}

\titleformat{\subsubsection}
  {\normalfont\normalsize\itshape}
  {}{0pt}{}
\titlespacing*{\subsubsection}{0pt}{10pt}{4pt}

% 章题下的俄文原题
\newcommand{\chaptersubtitle}[2]{%
  \begin{center}
  \vspace{-14pt}
  {\small\sffamily\color{pencil}#1}\\[2pt]
  {\small\sffamily\itshape\color{pencil}#2}
  \end{center}
  \vspace{4pt}
}

% 章前题记
\newcommand{\epigraphbox}[2]{%
  \begingroup
  \vspace{2pt}
  \begin{list}{}{\leftmargin=3.2em \rightmargin=2.2em \topsep=0pt
                 \parsep=0pt \itemsep=0pt}
  \item\relax
  {\small\linespread{1.32}\selectfont\color{ink!88}#1\par}
  \vspace{3pt}
  {\small\raggedleft\color{pencil}#2\par}
  \end{list}
  \vspace{6pt}
  \endgroup
}

% ---------------------------------------------------------------- 页眉页脚
% 注意：fancyhdr v5 的 \fancypagestyle 在使用时会执行其中的 \fancyhf{}，
% 从而清空全局 fancy 设置。故此处直接定义 \ps@ 样式，绕开该机制。
\makeatletter
\def\ps@frontmatterpage{%
  \let\@mkboth\@gobbletwo
  \let\@oddhead\@empty \let\@evenhead\@empty
  \def\@oddfoot{\hfil\small\thepage\hfil}%
  \def\@evenfoot{\hfil\small\thepage\hfil}}
\def\ps@plain{%
  \let\@mkboth\@gobbletwo
  \let\@oddhead\@empty \let\@evenhead\@empty
  \def\@oddfoot{\hfil\small\thepage\hfil}%
  \def\@evenfoot{\hfil\small\thepage\hfil}}
\makeatother

\pagestyle{fancy}
\fancyhf{}
\renewcommand{\headrulewidth}{0.3pt}
\renewcommand{\headrule}{\hbox to\headwidth{\color{rulegrey}\leaders\hrule height 0.3pt\hfill}}
\fancyhead[LE]{\small\color{pencil}\itshape 战争的雷霆}
\fancyhead[RO]{\small\color{pencil}\itshape\leftmark}
\fancyfoot[C]{\small\color{ink}\thepage}

\renewcommand{\chaptermark}[1]{\markboth{第\,\thechapter\,章\quad #1}{第\,\thechapter\,章\quad #1}}
\renewcommand{\sectionmark}[1]{}

% ---------------------------------------------------------------- 脚注
\renewcommand{\footnoterule}{%
  \kern-3pt
  \hrule width 3.1cm height 0.4pt\kern 2.6pt}
\setlength{\skip\footins}{16pt}
\renewcommand{\footnotesize}{\fontsize{8.1pt}{11.4pt}\selectfont}
\makeatletter
\renewcommand{\@makefnmark}{\hbox{\normalfont\scriptsize\@thefnmark}}
\makeatother
\renewcommand{\thefootnote}{\arabic{footnote}}
\renewcommand{\footnotelayout}{\linespread{1.18}\selectfont}

% ---------------------------------------------------------------- 图表
\captionsetup{
  font={small}, labelfont={bf,sf}, labelsep=quad,
  justification=centering, skip=7pt
}
\captionsetup[table]{position=top, skip=5pt}
\captionsetup[figure]{position=bottom, skip=7pt}

\newcommand{\figsource}[1]{%
  \par\vspace{2pt}{\footnotesize\raggedright\color{pencil}资料来源：#1\par}}
\newcommand{\tabsource}[1]{%
  \par\vspace{1pt}{\footnotesize\raggedright\color{pencil}资料来源：#1\par}}

% ---------------------------------------------------------------- 档案框
\newtcolorbox{archive}[2][]{%
  enhanced, breakable,
  colback=shade, colframe=graphite!80,
  boxrule=0.45pt, arc=0pt, outer arc=0pt,
  left=13pt, right=13pt, top=10pt, bottom=10pt,
  fontupper=\footnotesize\linespread{1.3}\selectfont,
  before skip=13pt, after skip=13pt,
  overlay={%
    \node[anchor=north west, inner sep=0pt, xshift=13pt, yshift=-6pt]
      at (frame.north west)
      {\sffamily\scriptsize\color{graphite}\MakeUppercase{#2}};
    \node[anchor=north east, inner sep=0pt, xshift=-13pt, yshift=-4pt]
      at (frame.north east)
      {\sffamily\scriptsize\color{stamp}\bfseries #1};
    \draw[stamp, line width=0.4pt]
      ([xshift=-13pt,yshift=-8pt]frame.north east) -- ++(0,-5pt);
  },
  before upper={\vspace{13pt}},
}

\newcommand{\archivestamp}[1]{}

% 小号“译者注”标记
\newcommand{\yizhe}[1]{{\footnotesize\color{pencil}［译者注：#1］}}

% 强调
\newcommand{\kw}[1]{\textbf{#1}}
\newcommand{\ru}[1]{{\itshape #1}}

% 引用文献的斜体书名（中文用书名号，不另设）
\newcommand{\citebk}[1]{《#1》}

% TOC: 使用 book 类默认设置
\makeatletter
\renewcommand{\@pnumwidth}{2.4em}
\renewcommand{\@tocrmarg}{3.4em}
\renewcommand{\@dotsep}{1.1}
\makeatother

% 未编号条目（附录、序言等）
\newcommand{\tocchap}[1]{%
  \addcontentsline{toc}{chapter}{\texorpdfstring{#1}{#1}}}

% ---------------------------------------------------------------- 杂项
\newcommand{\blankpage}{\null\thispagestyle{empty}\newpage}
\setlength{\parskip}{0pt}

% 每页页眉下细线通过 header 中的 \headrule 关闭，改用文字规则
\newcommand{\thinrule}{\par\vspace{2pt}{\color{rulegrey}\hrule height 0.3pt}\vspace{2pt}}

% 文本内的“印章”
\newcommand{\seal}[2][stamp]{%
  \begin{tikzpicture}[baseline=-0.6em]
    \node[draw=#1, line width=0.6pt, text=#1, inner sep=2.4pt,
          font=\sffamily\scriptsize, rotate=-6, fill=white]
      {#2};
  \end{tikzpicture}}

\endinput
